% D1 — maquette du référentiel. Contenu : bir/referentiel.md.
\documentclass{BIRdoc}
\BIRdocument{Référentiel}{Brevet d'Initiation à la Radio}
\newcommand{\dom}[2]{\section*{Domaine #1 --- #2}}
\newenvironment{competences}{%
  \small\tabularx{\linewidth}{P{33mm}LLP{20mm}}
  \BIRentete \BIRth{Séance} & \BIRth{L'élève sait…} & \BIRth{L'élève sait faire…} & \BIRth{Certificat} \\}
  {\bottomrule\endtabularx\par}
\begin{document}

\BIRtitre{Référentiel du Brevet d'Initiation à la Radio}{Ce que l'élève sait et sait faire à l'issue des 24 séances}

\section*{Objet}

Le Brevet d'Initiation à la Radio (BIR) atteste qu'un élève de collège ou de
lycée connaît les bases de la radio et sait les mettre en œuvre : écouter,
mesurer, construire une antenne simple, tenir sa place dans un contact.

Il se prépare en 24 séances de 1 h 30, par groupes réduits, en présence d'un
radioamateur. Il ne donne aucun droit d'émettre : ce droit s'obtient par le
certificat d'opérateur des services d'amateur, auquel le BIR prépare le
terrain.

\section*{Cinq domaines}

\begin{tabularx}{\linewidth}{P{18mm}LP{45mm}}
\BIRentete \BIRth{Domaine} & \BIRth{Intitulé} & \BIRth{Séances} \\
A & Découvrir la radio et les radioamateurs & 1, 20 \\
B & Électricité et électronique & 4, 5, 6, 17, 18, 19 \\
C & Ondes, spectre et propagation & 7, 8, 9, 10 \\
D & Modulations, antennes et station & 11 à 16 \\
E & Trafic, règles et sécurité & 2, 3, 21, 22, 23 \\
\bottomrule
\end{tabularx}

\begin{BIRencadre}{Lire la colonne « certificat »}
Elle rattache chaque séance au programme du certificat d'opérateur (arrêté
du 21 septembre 2000, annexe I) : \textbf{R} pour la première partie,
réglementation, \textbf{T} pour la seconde, technique.
\end{BIRencadre}

\dom{A}{Découvrir la radio et les radioamateurs}
\begin{competences}
1 · La radio autour de nous & citer cinq usages de la radio ; dire ce qu'est un radioamateur et ce qui le distingue d'un auditeur & régler un récepteur sur une station et la décrire & R ch. 1, 2 \\
20 · La radio en action & expliquer le rôle d'un relais et d'un satellite ; dire ce qu'est un locator & trouver son locator ; localiser un émetteur caché & R ch. 2 \\
\end{competences}

\dom{B}{Électricité et électronique}
\begin{competences}
4 · Tension et courant & définir tension et courant, avec leurs unités ; distinguer conducteur et isolant ; citer les dangers du courant & câbler un circuit simple ; mesurer une tension et un courant & T 1.1, 1.2 · R 3.5 \\
5 · La résistance et la loi d'Ohm & énoncer la loi d'Ohm ; lire le code des couleurs ; reconnaître les symboles usuels & calculer une des trois grandeurs à partir des deux autres ; vérifier par la mesure & T 2.1, 3.1, 7 \\
6 · Puissance et énergie & relier puissance, tension et courant ; dire ce que mesure une capacité de batterie & calculer une puissance et une autonomie & T 1.9, 3.3 \\
17 · Le récepteur & nommer les étages d'un récepteur et leur rôle ; dire à quoi sert le squelch & régler fréquence, volume, mode et squelch d'un poste & T 4 \\
18 · L'émetteur & nommer les étages d'un émetteur ; dire à quoi sert une charge fictive ; citer une cause de perturbation & lire une puissance sur un wattmètre & T 5 · R 3 \\
19 · L'atelier de soudure & reconnaître résistance, condensateur, diode, transistor & réaliser une soudure propre ; suivre un schéma d'implantation & T 2 \\
\end{competences}

\dom{C}{Ondes, spectre et propagation}
\begin{competences}
7 · Du courant alternatif à l'onde & définir fréquence, période et longueur d'onde ; relier fréquence et longueur d'onde & calculer une longueur d'onde ; lire une période sur un oscillogramme & T 1.5, 1.6 \\
8 · Le spectre et les bandes & situer HF, VHF et UHF ; dire pourquoi les fréquences sont partagées et par qui ; citer six bandes radioamateur & lire un spectre et une chute d'eau ; dresser une carte du spectre & R ch. 1, 2 \\
9 · La propagation à vue & expliquer l'horizon radio et l'effet du relief & mesurer une portée et la reporter sur un plan & T 6.1 \\
10 · Les ondes courtes et l'ionosphère & expliquer comment une onde courte franchit les océans ; dire ce qui change entre le jour et la nuit & identifier une station lointaine et la placer sur une carte & T 6.1 \\
\end{competences}

\dom{D}{Modulations, antennes et station}
\begin{competences}
11 · Moduler & dire ce que moduler veut dire ; distinguer AM, FM et SSB ; définir la largeur de bande & reconnaître un mode à l'oreille et à l'écran & T 1.8 \\
12 · Le Morse & expliquer le principe du code ; connaître SOS et les lettres de son prénom & émettre et recevoir un mot court & T 1.8 \\
13 · Les modes numériques et l'image & distinguer analogique et numérique ; citer deux modes numériques & décoder une transmission avec un logiciel & T 1.8, 1.10 \\
14 · Les antennes : principes & dire à quoi sert une antenne ; relier la longueur d'un dipôle à la longueur d'onde ; définir la polarisation & calculer la longueur d'un dipôle & R 4.1, 4.2 · T 6.2 \\
15 · Construire son antenne & distinguer antenne omnidirectionnelle et antenne directive & assembler une antenne ; constater sa directivité & R 4.1 \\
16 · Câbles, connecteurs et réglage & décrire un câble coaxial ; dire ce qu'indique le rapport d'ondes stationnaires & mesurer et améliorer le réglage d'une antenne & R 4.3 · T 6.3, 7 \\
\end{competences}

\dom{E}{Trafic, règles et sécurité}
\begin{competences}
2 · Se présenter à l'antenne & lire un indicatif ; connaître l'alphabet d'épellation & épeler et noter un indicatif sans erreur & R ch. 2, 5 \\
3 · Un contact radio & décrire les étapes d'un contact ; donner un report ; connaître dix codes Q & tenir un carnet de trafic & R ch. 5 \\
21 · Préparer son contact & dire ce qu'on annonce, et dans quel ordre & dérouler un contact complet, script en main & R ch. 5 \\
22 · Au micro & --- & tenir sa place dans un vrai contact, aux côtés du radioamateur & --- \\
23 · Les règles, et après ? & dire qui a le droit d'émettre ; décrire l'examen du certificat d'opérateur ; nommer l'ANFR & --- & R ch. 1, 2 \\
\end{competences}

\section*{L'examen}

L'examen se passe en ligne. La plateforme tire les questions dans la banque
du BIR et mélange les réponses.

\begin{tabularx}{\linewidth}{P{30mm}L}
\BIRentete \BIRth{} & \BIRth{Proposition} \\
Forme & questionnaire à choix unique en ligne, quatre réponses proposées \\
Durée & 1 heure \\
Questions & 40, soit 8 par domaine \\
Réussite & 20 sur 40, avec au moins 3 bonnes réponses dans chaque domaine \\
Savoir-faire & validés en cours de formation dans le livret de l'élève, sans note \\
\bottomrule
\end{tabularx}

\section*{Trois niveaux d'exigence}

\begin{enumerate}
\item \textbf{Reconnaître} --- nommer, situer, associer.
\item \textbf{Expliquer} --- dire pourquoi, avec ses mots.
\item \textbf{Mettre en œuvre} --- calculer, mesurer, régler.
\end{enumerate}

Le BIR vise les niveaux 1 et 2 partout, et le niveau 3 sur les savoir-faire
des fiches d'activité.

\end{document}
